\documentclass[11pt]{article}
\usepackage{setspace}
\usepackage{geometry}
\geometry{margin=1in}
\setstretch{1.12}

\begin{document}

% --- Applicant Information ---
\noindent
\textbf{Zhongli Wang} \\
+49 160 9541 6015 \\
nkwangzhongli@gmail.com \\
https://sites.google.com/view/zhongli-wang \\
Bielefeld, Germany \\

\vspace{0.5cm}

% --- Recipient Information ---
\noindent
\textbf{Prof.\ Sascha Steffen} \\
Centre for European Transformation \\
Frankfurt School of Finance \& Management \\
Frankfurt am Main, Germany \\

\vspace{0.5cm}

% --- Subject Line ---
\noindent
\textbf{Application for Postdoctoral Fellowship in Finance --- Centre for European Transformation} \\

\vspace{0.3cm}

\noindent Dear Prof.\ Steffen and Members of the Selection Committee,\\[0.2cm]

I am writing to apply for the Postdoctoral Fellowship in Finance at the Centre for European Transformation. 
I am currently completing my Ph.D.\ in Economics at Bielefeld University and Universit\`a Cattolica del Sacro Cuore within the Marie Sk{\l}odowska--Curie EPOC programme. 
My research combines applied microeconomics, financial economics, and computational methods to study how market structure and information frictions shape firms' access to funding and their real decisions. 
I am particularly interested in the frictions that hinder firms from borrowing, rolling over debt, or raising equity---and in how policy design and market architecture can strengthen resilience during stress episodes such as the Global Financial Crisis and the March 2020 dash-for-cash.

A core theme of my doctoral work is the interaction between constraints (capacity, liquidity, regulation, and information) and dynamic adjustment by heterogeneous agents. 
My dissertation develops quantitative models of technology-intensive markets in which firms choose pricing and investment under network effects, congestion, and imperfect information. 
Methodologically, I combine structural modeling and simulation-based estimation with scalable computation (Python, parallel simulation, Monte Carlo methods) to quantify behavioral mechanisms and evaluate counterfactual scenarios. 
This training transfers directly to research on financial intermediation and debt markets, where rollover risk, liquidity provision, and institutional rules interact with heterogeneous balance sheets to shape investment and growth.

In parallel, I have pursued research on financial market behavior and transparency. 
In \textit{Polytope Fraud Theory} (IRFA, 2025), I use supervised machine learning to detect structured irregularities in firm financial statements, focusing on firms targeted by activist short sellers in U.S.\ markets. 
The project links institutional incentives and monitoring frictions to observable anomalies and proposes a framework for investor protection and market integrity. 
More broadly, it sharpened my experience with firm-level data, balance-sheet variables, and model evaluation under noisy measurement---skills directly relevant for studying financing constraints, credit allocation, and the distributional consequences of funding access.

A second coauthored project, \textit{Agentic Workflows for Economic Research}, develops a structured framework for integrating large-language-model (LLM) tools into empirical workflows to improve transparency, reproducibility, and efficiency. 
This work is directly relevant to the Centre's interest in data science and unstructured data: it provides practical building blocks for collecting and structuring text-based evidence (e.g., earnings calls, regulatory documents, news) and for creating auditable pipelines that link unstructured signals to firm outcomes, financing conditions, and policy changes. 
I am excited to apply these tools to capital-market questions where text and narrative components are economically central---for example, in measuring funding conditions, expectations, and policy transmission during crises.

At Frankfurt School and the Centre for European Transformation, I would aim to contribute to research on how market structure, contracting, and policy jointly shape investment incentives and resilience in Europe. 
In particular, I am interested in projects on (i) debt rollover and refinancing risk under stress, (ii) the design and effectiveness of long-term contracting and public guarantees in financing innovation and the energy transition, and (iii) fragmentation and access to capital across countries, firm size classes, and sectors. 
My comparative advantage is to combine theory-driven economic mechanisms with modern computation: I can contribute to model development, quantitative simulation, and empirical validation using firm-level data, and I can extend the research frontier by incorporating unstructured data (e.g., disclosures, earnings calls, policy documents) into measurement and identification strategies.

Beyond research, I would welcome the opportunity to contribute to the Centre's broader mission through policy-oriented outputs, presentations, and grant applications. 
My doctoral training involved coordinating computationally intensive projects across institutions and producing reproducible code and documentation, which aligns well with the Centre's emphasis on high-quality, team-based research and impact. 
I am enthusiastic about working in a collaborative environment that links frontier academic research to real-world policy debates on Europe's capital markets and long-run transformation.

Thank you so much for considering my application! 
I would be delighted to discuss how my research background and quantitative skill set can contribute to the Centre for European Transformation and its work on financial intermediation, debt capital markets, and the design of robust European capital markets.\\[2mm]

\noindent Sincerely,\\[12mm]
\noindent Zhongli Wang

\end{document}